\chapter{Environnement technique}
\label{ch:environnement}

\section{Vue d'ensemble de la pile technique}

\projet{} repose sur une pile volontairement légère, dont je maîtrise chaque brique.

\begin{table}[h]
\centering
\renewcommand{\arraystretch}{1.3}
\small
\begin{tabularx}{\textwidth}{@{}L{3.4cm} L{4.2cm} X@{}}
\toprule
\textbf{Couche} & \textbf{Technologie} & \textbf{Rôle} \\
\midrule
Langage & PHP 8.3 (\texttt{declare(strict\_types=1)}) & Logique applicative \\
SGBD & PostgreSQL 16 via PDO & Persistance relationnelle \\
Serveur applicatif & FrankenPHP (Caddy intégré) & Exécution de PHP, service du statique \\
Reverse proxy / edge & Traefik v3 & Routage, terminaison TLS (Let's Encrypt) \\
Conteneurisation & Docker, \texttt{docker compose} & Reproductibilité dev = prod \\
Dépendances & Composer & \texttt{minify}, \texttt{phpdotenv}, \texttt{psr/log} \\
Front-end & HTML, CSS, JavaScript vanilla & Interfaces et appels AJAX (\texttt{fetch}) \\
Gestion de versions & Git & Historique et traçabilité \\
\bottomrule
\end{tabularx}
\caption{Pile technique de \projet{}.}
\label{tab:stack}
\end{table}

\section{Justification des choix techniques}

Chaque choix est un compromis explicite, formulé selon le schéma « X plutôt que Y parce que Z,
en acceptant le coût W ».

\begin{description}
  \item[MVC maison plutôt qu'un framework (Symfony, Laravel).] Objectif pédagogique :
        \emph{comprendre et posséder} le routeur, l'autoloader et le cycle requête / réponse,
        plutôt que de déléguer à de la « magie » non maîtrisée. Coût assumé : réécrire des
        briques qu'un framework fournirait. En contexte professionnel, un framework serait
        retenu.
  \item[PostgreSQL plutôt que MariaDB.] Typage strict, séquences (\texttt{SERIAL}), robustesse
        et conformité SQL. Le projet a d'ailleurs été \emph{migré} de MariaDB vers PostgreSQL,
        ce qui a imposé d'adapter le schéma (\texttt{AUTO\_INCREMENT} $\to$ \texttt{SERIAL},
        \texttt{datetime} $\to$ \texttt{TIMESTAMP}).
  \item[FrankenPHP plutôt qu'Apache ou Nginx + PHP-FPM.] Un seul binaire embarquant Caddy,
        terminaison TLS automatique, image Docker simple.
  \item[Traefik plutôt que Nginx en reverse proxy.] Configuration \emph{par labels Docker},
        certificats Let's Encrypt automatiques, découverte dynamique des services.
  \item[PDO plutôt qu'une extension native.] Abstraction portable et, surtout, \textbf{requêtes
        préparées} qui neutralisent l'injection SQL (section~\ref{sec:acces}).
\end{description}

\section{Poste de travail et outils}

L'environnement de développement réunit : un IDE, \textbf{Git} pour la gestion de versions
(historique de \emph{commits} conventionnels, une fonctionnalité par série de commits),
\textbf{Composer} pour les dépendances (\file{composer.json} / \file{composer.lock}), et une
gestion rigoureuse des \textbf{secrets}. Les identifiants de base de données ne figurent jamais
dans le code : ils sont chargés depuis un fichier \file{.env} (ignoré par Git) au moyen de la
bibliothèque \texttt{vlucas/phpdotenv}, un gabarit \file{.env.example} documentant les variables
attendues.

\begin{lstlisting}[style=php, caption={Chargement des secrets et connexion PDO — \file{includes/database.php}}]
$envPath = dirname(__DIR__);
if (file_exists($envPath . '/.env')) {
    $dotenv = Dotenv\Dotenv::createImmutable($envPath);
    $dotenv->load();
}
// ... lecture de DB_HOST, DB_NAME, DB_USER, DB_PASSWORD ...
$dsn = "pgsql:host={$host};port={$port};dbname={$dbname};options='--client_encoding=UTF8'";
$options = [
    PDO::ATTR_ERRMODE            => PDO::ERRMODE_EXCEPTION,
    PDO::ATTR_DEFAULT_FETCH_MODE => PDO::FETCH_ASSOC,
    PDO::ATTR_EMULATE_PREPARES   => false,   // vraies requêtes préparées côté serveur
];
$this->connection = new PDO($dsn, $user, $password, $options);
\end{lstlisting}

La désactivation de \texttt{ATTR\_EMULATE\_PREPARES} force de \emph{vraies} requêtes préparées
côté serveur PostgreSQL, renforçant la protection contre l'injection SQL.

\section{Conteneurisation (Docker)}

L'application est décrite par deux piles \texttt{docker compose} : l'une pour le développement,
l'autre pour la production.

\paragraph{Développement.} Traefik en HTTP, FrankenPHP en \emph{bind-mount} du code (rechargement
à chaud), PostgreSQL exposé sur le port 5432 (pour les outils locaux), et un service
\texttt{migrate} exécuté une seule fois avant le démarrage de l'application. Un \emph{healthcheck}
PostgreSQL garantit l'ordre de démarrage.

\begin{lstlisting}[style=yaml, caption={Extrait du \file{docker-compose.yml} (développement) : service applicatif et base}]
  app:
    build: { context: ., dockerfile: Dockerfile }
    image: solage-app
    volumes:
      - ./:/app          # rechargement du code à chaud en dev
      - /app/vendor      # vendor/ provient de l'image
    depends_on:
      postgres: { condition: service_healthy }
      migrate:  { condition: service_completed_successfully }
    labels:
      - traefik.enable=true
      - traefik.http.routers.solage.rule=Host(`localhost`)

  postgres:
    image: postgres:16-alpine
    healthcheck:
      test: ["CMD-SHELL", "pg_isready -U $${POSTGRES_USER} -d $${POSTGRES_DB}"]
      interval: 5s
\end{lstlisting}

\paragraph{Production.} Traefik en HTTPS (Let's Encrypt, redirection HTTP $\to$ HTTPS), en-tête
HSTS, FrankenPHP servi depuis l'image (sans \emph{bind-mount}), et PostgreSQL \emph{non exposé}
(accessible uniquement sur le réseau interne). Le détail figure en
section~\ref{sec:deploiement} (préparation du déploiement).

